\documentclass{llncs}
%\usepackage{fullpage}
%\usepackage{slashbox}
\usepackage{relsize}
\usepackage{amsmath}
\usepackage{amssymb}
\usepackage{oz}
%\usepackage{lmodern}
\usepackage{listings} 
%\usepackage{fixmath}
\usepackage{csp-oz}
\usepackage{delarray}
\usepackage{slashed}
%\usepackage{graphicx}
\usepackage{etex}
\usepackage{pgf}
\usepackage{tikz}
\usepackage{dsfont} 
%\usepackage{arevmath}
\usepackage{float} 
%\usepackage{algorithm2e}
%\usepackage[noend]{algorithm2e}
%\usepackage[noend,linesnumbered,ruled,vlined]{algorithm2e}
\usepackage[noend,linesnumbered,ruled,vlined]{algorithm2e}
%\usepackage[noend]{algorithmic}
%\usepackage{times}
\usetikzlibrary{fadings}
\usetikzlibrary{positioning,automata}
\usepackage{setspace}
%\usepackage{bbm}
%\usepackage{dsfont}
\usetikzlibrary{arrows,automata,shapes,patterns}
\usepackage{pgfplots}
\usepackage{filecontents}
%\usepackage{bbding}
\usepackage{pifont}% http://ctan.org/pkg/pifont
\newcommand{\cmark}{\ding{51}}%
\newcommand{\xmark}{\ding{55}}%
%\usepackage{bbold}
%\usepackage{bbding}
%\usetikzlibrary{decorations}
%\usetikzlibrary{snakes}
\usepackage{stackengine}
\usetikzlibrary{matrix}
\usepackage{multirow}
%\usepackage{turnstile}
\def\ruleoffset{1pt}
\newcommand\specialvdash[2]{\mathrel{\ensurestackMath{
  \mkern2mu\rule[-\dp\strutbox]{.4pt}{\baselineskip}\stackon[\ruleoffset]{
    \stackunder[\dimexpr\ruleoffset-.5\ht\strutbox+.5\dp\strutbox]{
      \rule[\dimexpr.5\ht\strutbox-.5\dp\strutbox]{2.5ex}{.4pt}}{
        \scriptstyle #1}}{\scriptstyle#2}\mkern2mu}}
}
\DeclareRobustCommand\model{\mathrel{|}\joinrel\mkern-.5mu\mathrel{-}}

\DeclareRobustCommand{\invE}{\text{\reflectbox{\reflectbox{$\exists$}}}}

    \newcommand{\invertleadsto}{%
    \mathrel{\reflectbox{\rotatebox[origin=c]{180}{$E$}}}}

\SetAlFnt{\footnotesize}
\SetAlCapFnt{\footnotesize}
\SetAlCapNameFnt{\footnotesize}

\newcommand{\coop}[2][]{\langle\!\langle{#2}\rangle\!\rangle_{_{\!\mathit{#1}}}\,}

\SetKwFor{For}{for}{do}{EndLoop}
\SetKwFor{ForAllP}{loop forever do}{}{end}

\tikzset{
  treenode/.style = {align=center, inner sep=0pt, text centered,
    font=\sffamily},
  arn_o/.style = {thick,minimum height=1.3em,treenode, rectangle font=\sffamily\bfseries, draw=black,
    fill=orange, text width=1.3em},% arbre rouge noir, noeud noir
     arn_bo/.style = {thick,minimum height=2.5em,treenode, rectangle font=\sffamily\bfseries, draw=black,
    fill=orange, text width=2.5em},% arbre rouge noir, noeud noir
  arn_g/.style = {thick,minimum height=1.3em,treenode, rectangle, fill=mygreen, draw=black, 
    text width=1.3em },% arbre rouge noir, noeud rouge
      arn_so/.style = {thick,minimum height=0.75em,treenode, rectangle font=\sffamily\bfseries, draw=black,
    fill=orange, text width=0.75em},% arbre rouge noir, noeud noir
  arn_sg/.style = {thick,minimum height=0.75em,treenode, rectangle, fill=mygreen, draw=black, 
    text width=0.75em },% arbre rouge noir, noeud rouge
  arn_xg/.style = {fill=mygreen, draw=black,thick,treenode, rectangle, draw=black,
    minimum width=0.75em, minimum height=0.75em},% arbre rouge noir, nil
     arn_xo/.style = {fill=orange, draw=black,thick,treenode, rectangle, draw=black,
    minimum width=1em, minimum height=1em},% arbre rouge noir, nil
}



\usepackage{url}
\urldef{\mails}\path|u18074848@tuks.co.za|
\urldef{\mailsb}\path|ntimm@cs.up.ac.za|    


\usepackage{ntheorem}
%\newtheorem{theorem}{Theorem}
\newtheorem*{theorem2}{Theorem 2.}
%\newtheorem{theorem2}{Theorem 2.}
%\newtheorem*{theorem2}{Theorem}
\newtheorem*{theorem3}{Lemma 1.}
\newtheorem*{theorem4}{Lemma 2.}

\newcommand{\BIGOP}[1]
{
\mathop{\mathchoice%
{\raise-0.22em\hbox{\huge $#1$}}%
{\raise-0.05em\hbox{\Large $#1$}}{\hbox{\large $#1$}}{#1}}}
\newcommand{\bigtimes}{\BIGOP{\times}}
% nur fuer Bigboxplus andere Korrekturen
\newcommand{\BIGboxplus}{\mathop{\mathchoice%
{\raise-0.35em\hbox{\huge $\boxplus$}}%
{\raise-0.15em\hbox{\Large $\boxplus$}}{\hbox{\large $\boxplus$}}{\boxplus}}}

\newcommand{\subscriptedfrac}[3]{%
 \frac{\,#1\,}{\,#2\,}{}_{#3}%
}

\newcommand{\at}{$@$}

\newcommand\bigforall{\mbox{\Large $\mathsurround0pt\forall$}} 

\newenvironment{keywords}{
       \list{}{\advance\topsep by0.35cm\relax\small
       \leftmargin=0cm
       \labelwidth=0.35cm
       \listparindent=0.35cm
       \itemindent\listparindent
       \rightmargin\leftmargin}\item[\hskip\labelsep
                                     \bfseries Keywords:]}
     {\endlist}

\begin{document}
\title{Strategy Synthesis for Competitive Region Control via 
Q-Learning-Assisted Formal Reasoning 
} 

\author{
Steven Jordaan\orcidID{0009-0005-2928-2122}  \and
Nils Timm\orcidID{0000-0002-9656-3240} 
}
\institute{Department of Computer Science, University of Pretoria, Pretoria, South Africa \\
\mails \\ \mailsb
}

\maketitle
\begin{abstract}
\end{abstract}

\section{Introduction} 
Discussion of real-world mission/safety critical multi-agent problems and the need strategies that guarantee mission success.
\section{Related Work}


\section{Region Control Multi-Agent Systems}
In our approach we introduce  \emph{region control multi-agent systems} (RCMAS).
\begin{definition}[Region Control Multi-Agent Systems] 
A region control multi-agent system is a tuple $M = (Agt, T, o)$ where
\begin{itemize}
 \item $Agt = \{a_k \mid 1 \leq k \leq n \}$ is a finite set of agents,
\item $T \subseteq  \{(i,j) \mid 1 \leq i \leq w,  1 \leq j \leq r,\}$ is a territory of sectors,
\item $o : Agt \rightarrow \mathbb{N}$ is an objective function that defines minimum size of a cohesive region that each agent needs to control.
\end{itemize}    
\end{definition}
We define $T$ as a \emph{subset} relation in order to exclude sectors from the territory that we regard as \emph{inaccessible} in RCMAS instances.  
Each agent $a$ has the qualitative mission of controlling a cohesive region at least size $o(a)$ by the time all sectors of the territory are occupied.  
Moreover, each agent has the quantitative mission of maximising the size of its largest cohesive region.
The actions that can be performed for this are as follows.

Todo: Formally define a cohesive region $R \subseteq T \ldots$ (see Definition \ref{def:cohesive-regions}). 
    
\begin{definition}[Actions]
Given an RCMAS $M$, the {set of actions}is $Act = \{ act^a_{i,j} \mid a \in Agt,  G(i,j)\}$.
\end{definition}

\noindent
Hence,  actions allow agents to take control of units of the territory.
This gives rise to the notion of \emph{states} of an RCMAS, which we subsequently define. 
An RCMAS runs in discrete rounds where in each round each agent chooses its next action.  In a round the tuple of chosen actions, one per agent, gets executed simultaneously.  The execution of actions leads to an evolution of the system between different states over time. 

\begin{definition}[States]
A state of an RCMAS $M$ is a function $s : T \rightarrow Agt^+$ where $Agt^+ = Agt \cup \{a_0\}$ and $a_0$ is a dummy agent.
If $s(i,j) = a_0$ then sector $r$ is unoccupied in state $s$. 
If $s(i,j) = a_k$ and $k > 0$ then $(i,j)$ is occupied by agent $a_k$ in $s$.
We denote by $s_0$ the initial state of $M$, where $s(i,j) = a_0$ for each $(i,j) \in T$, i.e., initially all sectors are unoccupied.
If we want to express that a sector $(i,j)$ is occupied by agent $a_k$ but the current state is not further specified, then we simply write $(i,j) \leftarrow a_k$.
We denote by $S$ the set of all possible states of $M$,  
and we implicitly assume that $S$ comprises the special state $s^x$,
that indicates mission failure. 
\end{definition}

\noindent
Hence, states describe which sectors are currently under the control of which agents.

\begin{definition}[Cohesive Regions]
\label{def:cohesive-regions}
Let $M$ be an RCMAS, $a \in Agt$, $s \in S$, and $s(i,j) = a$.
Then the cohesive region $R$ associated with sector $(i,j)$ 
is inductively defined as follows:
\begin{itemize}
\item $(i,j) \in R$;
\item if $(i',j') \in R$ and $s(i'+1,j') = a$, then $(i'+1,j') \in R$;
\item if $(i',j') \in R$ and $s(i'-1,j') = a$, then $(i'-1,j') \in R$;
\item if $(i',j') \in R$ and $s(i',j'+1) = a$, then $(i',j'+1) \in R$;
\item if $(i',j') \in R$ and $s(i',j'-1) = a$, then $(i',j'-1) \in R$;
\end{itemize}
\end{definition}



In each state only a subset of actions may be available for execution by an agent, which we define by the protocol:

\begin{definition}[Action Availability Protocol]
  The {action availability protocol} is a function $P : S \times Agt \rightarrow 2^{Act}$ defined for each $s \in S$ and $a \in Agt$:
  \[
  P(s,a) = \{ act^a_{i,j} \mid s(i,j) = a_0\}
  \]
\end{definition}
Thus,  an agent can attempt to occupy a sector that is currently unoccupied.
  
  \begin{definition}[Action Profiles]
  An action profile in an RCMAS $M$ is a mapping $ap : Agt \rightarrow Act$.  $AP$ denotes the set of all action profiles. We say that a profile $ap$ is executable in a state $s \in S$ if for each $a \in Agt$ we have that $ap(a) \in P(s,a)$.
  \end{definition}
 Based on action profiles we can formally define the evolution of an RCMAS.
  
   \begin{definition}[Evolution]
The evolution of an RCMAS is a relation $\delta \subseteq S \times AP \times S$ where $(s, ap, s') \in \delta$ iff $ap$ is executable in $s$ and for each $(i,j) \in T$:
\begin{enumerate}
\item if $s(i,j) = a_0$ then: %\\
\begin{enumerate}
\item if $\neg \exists a : ap(a) = act^a_{i,j}$ then $s'(i,j) = a_0$;
\item else if $\exists  a : ap(a) = act^a_{i,j} \wedge \forall a' \neq a : ap(a') \neq act^{a'}_{i,j}$ then $s'(i,j) = a$; %\\ 
\item otherwise $s' = s^x$; %\\
\end{enumerate}
\item if $s(i,j) = a$ for some $a \in Agt$ then $s'(i,j) = a$. %\\
  \end{enumerate}
  \end{definition}
  If an action profile is executed in a state of an RCMAS $M$,  this leads to a transition of $M$ into a corresponding successor state, i.e.,
    a change in the control of sectors according to the actions chosen by the agents. According to the evolution, the attempt to occupy of a currently unoccupied sector $(i,j)$ by an agent $a$ will only be successful if $a$ is the only agent that attempts to occupy $(i,j)$ in the current round. If multiple agents attempt to occupy the same sector at the same time, this will result in a collision an the immediate failure of the mission. 
Moreover, any occupancy of sectors by agents will persist until the system terminates. 

\section{Strategies}
We are interested in solving strategy synthesis problems with regard to region control multi-agent systems.
A strategy is defined as follows:
  
         \begin{definition}[Strategy]
     A strategy  of an agent $a \in Agt$ in an RCMAS is an injective function $\alpha_a : S \rightarrow Act$. 
A strategy can also be denoted by a relation $\alpha_a \subseteq S \times Act$ where $\alpha_a(s,act^a_{i,j}) = true$ iff $\alpha_a(s) = act^a_{i,j}$. 
A joint strategy is a tuple of strategies $\alpha_{Agt} = (\alpha_a)_{a \in Agt}$, one for each $a \in Agt$.
  \end{definition}
  \noindent
A strategy determines which action an agent will choose  in which  state.  

\textbf{Comment: We can still consider uniform strategies in our scenario, if we assume that agents can only partially observe states. 
Then agents may have the same state observation in different states.}

  The outcome of a strategy $\alpha_{Agt}$ in a state $s$ is a path.
  
   \begin{definition}[Outcome of a Strategy]
   Let $M$ be an MRA and $s$ the initial state of $M$. 
   Moreover, let $\alpha_{Agt}$ be a joint strategy of $Agt$. Then the outcome of $\alpha_{Agt}$ in state $s$ is a path
  \[
  \begin{array}{lllll}
  \pi(s, \alpha_{Agt}) & = & & s_0 \ldots s_m \mid s_0 = s \ \wedge \ \forall  0 \leq t \leq s_m : \vspace*{1.5mm}\ \\
  & & & \bigg(s_t,\big(\alpha_{a_1}(s_t), \ldots, \alpha_{a_n}(s_t)\big),s_{t+1}\bigg) \in \delta \vspace*{1.5mm}\ \\ 
  & & & \wedge \ s_m = s^x \ \vee \ \forall (i,j) \in T : s_m(i,j) \neq a_0.
  \end{array}
  \]
   \end{definition}
   Hence, the outcome of a strategy is a path that either ends in the mission failure state or in a state where all sectors of the territory are occupied.  

\textbf{Todo: Formalise mission objectives/goals, winning/optimal/equilibrium strategies.
}

\textbf{Todo: Consider extensions where agents have preferences in terms of sectors, and where sectors are associated with numerical rewards.
}



\bibliographystyle{splncs03}
\bibliography{references}




\end{document} 
